\section{Limitations and Scope}
The plant, optics and sensing remain assumed rather than calibrated against a receiver and actuator. No hardware-in-the-loop, flight or industrial deployment is demonstrated. The six-model uncertainty set excludes continuous inertia/lag intervals; intermediate simulations do not enlarge the certified set. Source dynamics are a held-input RK4 sampled model with unrounded arithmetic bounds, and the theorem does not cover intersample threshold crossings.

The sensor event requires correct marginal tails, scales and zero bias for the full declared horizon. It cannot be reused indefinitely across multiple missions without appropriate risk accounting. Fixed offset and divergence and the specified gain-decline bound must remain valid throughout the probe. A supervisor cannot infer all these hypotheses from the same noisy data it is using to diagnose an unknown fault. Unobserved violations can escape monitoring.

The evaluation is development-guided and deliberately concentrated on admission boundaries. The baseline extension is frozen after the original outcomes, equal time does not equal equal excitation, and its adapted implementations do not constitute faithful reproductions of all close published methods. Full-text access and system-level mapping remain necessary for the closest recent optical work. No independent external mathematical review or exhaustive novelty search has been completed. These limitations preclude a claim of established Transactions novelty, physical safety, flight readiness or universal superiority.

Source attribution and content-assistance records are supplied. The repository's API reports no license; attribution alone does not establish redistribution rights. Human source-rights confirmation and approval of the author list, original contribution, disclosure and final content remain submission prerequisites. No author names, institutional affiliations or approvals are invented.

\section{Conclusion}
A reference-relative optical bound, simultaneous sensor event and finite-model actuator enclosure provide a conditional sampled probe-authorization rule. The study demonstrates additional diagnostic access in one restricted region, conservative refusal in the converse region, and substantial excitation/acceptance tradeoffs against adapted parameter filtering and tracking baselines. Zero observed outage does not validate a heuristic guard or a physical failure target. Approximate estimator confidence, hardware calibration, continuous actuator uncertainty and post-probe recovery remain separate unresolved problems.

\section*{Reproducibility}
The supplementary material supplies complete protocols, per-run rows, all new raw traces, independent calculation reports, symbolic identities, unit checks, stage-enclosure details, tests and source hashes. Earlier datasets and statistical corrections are preserved and identified separately; their records are not silently pooled with the 14,400 matched runs. The implementation and complete reproducibility package are identified by the supplied release manifest rather than an invented accepted-paper status.

\section*{Acknowledgments and AI-Use Disclosure}
OpenAI Codex assisted the mathematical formulation and checking, algorithm implementation, experimental scripts, analysis, literature screening, figures and drafting of Sections I--X and the supplementary material. MATLAB generated the recorded plant simulations; separate Python calculations supplied numerical replay and exact-algebra checks. This assistance is substantive and is disclosed under IEEE guidance \cite{aipolicy}. No independent human review is implied. Human authors must verify and take responsibility for the final scientific content, attribution and permissions before submission.
